\chapter{Appendix \thechapter}\label{appendixB}
\appendixchaptername{Experimental Data and Evaluation Results}

This appendix presents the quantitative datasets, evaluation 
tables, and representative examples that support the 
experimental analysis discussed in Chapter~4.  
The results include model selection metrics, human grading 
statistics, comparative analyses between AI and human 
evaluations, and detailed error case studies.

\section{Model Selection Dataset and WPS Calculation}

\rtab{tab:wps_result_extended} provides an extended 
version of the model selection results, including normalized 
accuracy, reasoning quality, and Weighted Performance Score 
(WPS) across all candidate LLMs. The WPS metric combines 
scoring precision, rubric alignment, and explanation clarity 
to select the most balanced model for deployment.

Each model was evaluated on a shared subset of 60 Thai 
written responses covering all Bloom’s taxonomy levels. 
Selection prioritized diversity in model architecture 
(decoder-only vs. instruction-tuned), accessibility 
(open vs. closed models), and regional adaptation 
(multilingual vs. Thai-optimized).

{
\setlength{\LTleft}{-2.2cm}
\setlength{\LTright}{0pt}
\centering
\begin{longtable}{l c l c l}
    \caption{Extended Model Selection Results with Normalized Scores and WPS}
    \label{tab:wps_result_extended}
    \\

    % For Header
    \hline
    \textbf{Model} & \textbf{Score Accuracy (\%)} & \textbf{Reasoning Quality} & \textbf{WPS (0–1)} & \textbf{Notes} \\
    \hline
    \endfirsthead

    % For Overflow Header
    \hline
    \textbf{Model} & \textbf{Score Accuracy (\%)} & \textbf{Reasoning Quality} & \textbf{WPS (0–1)} & \textbf{Notes} \\
    \hline
    \endhead

    % For Footer
    \hline
    \endfoot

    ChatGPT & 85.6 & High & 0.91 & Stable, high rubric alignment \\
    Gemini & 82.3 & High & 0.88 & Strong contextual reasoning \\
    Claude & 80.7 & High-Moderate & 0.84 & Excellent coherence, minor lexical bias \\
    Typhoon  & 77.2 & Moderate & 0.79 & Thai-adapted, concise explanations \\
    Grok & 71.5 & Moderate-Low & 0.72 & Frequent rubric omissions \\
    
\end{longtable}
}

%%%%%%%%%%%%%%%%%%%%%%%%%%%%%%%%%%%%%%%%%%%%%%%%%%%%%%%%%%%%

\section{Human Grading Data Summary}

\rtab{tab:teacher_summary} summarizes the distribution 
of scores assigned by the two expert raters (Teacher 1 and 
Teacher 2) under both non-rubric and rubric conditions.  
Overall agreement improved markedly when rubrics were 
introduced.

{
\setlength{\LTleft}{-1.35cm}
\setlength{\LTright}{0pt}
\centering
\begin{longtable}{l c c c l}
    \caption{Summary of Human Grading Results across Conditions}
    \label{tab:teacher_summary}
    \\

    % For Header
    \hline
    \textbf{Condition} & \textbf{Teacher} & \textbf{Exact Agreement (\%)} & \textbf{Average Score} & \textbf{Notes} \\
    \hline
    \endfirsthead

    % For Overflow Header
    \hline
    \textbf{Condition} & \textbf{Teacher} & \textbf{Exact Agreement (\%)} & \textbf{Average Score} & \textbf{Notes} \\
    \hline
    \endhead

    % For Footer
    \hline
    \endfoot

    Non-Rubric & Teacher 1 vs. 2 & 56.36 & 1.94 & Moderate consistency \\
    Rubric & Teacher 1 vs. 2 & 92.73 & 2.01 & Strong inter-rater alignment \\
    
\end{longtable}
}

\begin{figure}[H]
    \centering
    \includegraphics[width=0.7\textwidth]{agree}
    \caption{Confusion Matrix Teacher 1 compare to Teacher 2}
    \label{fig:confusion_matrices}
\end{figure}

%%%%%%%%%%%%%%%%%%%%%%%%%%%%%%%%%%%%%%%%%%%%%%%%%%%%%%%%%%%%

\section{AI vs. Human Score Comparison Tables}

An extended version of \rtab{tab:ai_human_comparison_app} 
lists per-Bloom-level Score Agreement Rate (SAR) and Cohen’s 
Kappa values. AI–human alignment was strongest in lower-order 
categories (Remembering, Understanding).

{
\setlength{\LTleft}{-0.5cm}
\setlength{\LTright}{0pt}
\centering
\begin{longtable}{l c c l}
    \caption{Detailed AI–Human Score Agreement by Bloom’s Category}
    \label{tab:ai_human_comparison_app}
    \\

    % For Header
    \hline
    \textbf{Bloom’s Level} & \textbf{SAR (\%)} & \textbf{Cohen's Kappa} & \textbf{Dominant Divergence Type} \\
    \hline
    \endfirsthead

    % For Overflow Header
    \hline
    \textbf{Bloom’s Level} & \textbf{SAR (\%)} & \textbf{Cohen's Kappa} & \textbf{Dominant Divergence Type} \\
    \hline
    \endhead

    % For Footer
    \hline
    \endfoot

    Remembering & 91.2 & 0.85 & Lexical synonym mismatch \\
    Understanding & 86.7 & 0.82 & Missing elaboration cues \\
    Applying & 81.3 & 0.77 & Ambiguous interpretation of examples \\
    Analyzing & 76.1 & 0.73 & Over-scoring detailed but incorrect logic \\
    Creating & 74.5 & 0.71 & Under-scoring creative, non-template answers \\
        
\end{longtable}
}

Representative divergences were also examined manually. 
For example, in the \textit{Creating} level, AI assigned a 
score of 2 to an original but unconventional response that 
both teachers rated 3, citing “insufficient keyword match.”

%%%%%%%%%%%%%%%%%%%%%%%%%%%%%%%%%%%%%%%%%%%%%%%%%%%%%%%%%%%%

\section{LLM Explanation Quality Examples}

A stratified subset of 200 AI-generated explanations was 
annotated by human evaluators for clarity, rubric alignment, 
and coverage of assessment criteria. 
\rtab{tab:explanation_quality} presents representative 
samples across quality levels.

{
\centering
\begin{longtable}{l p{0.8\textwidth}}
    \caption{Examples of AI Scoring Explanation Quality}
    \label{tab:explanation_quality}
    \\

    % For Header
    \hline
    \textbf{Quality Level} & \textbf{Example (Translated from Thai)} \\
    \hline
    \endfirsthead

    % For Overflow Header
    \hline
    \textbf{Quality Level} & \textbf{Example (Translated from Thai)} \\
    \hline
    \endhead

    % For Footer
    \hline
    \endfoot

    High & “The response explains both concepts correctly and connects them with appropriate terminology, matching the rubric’s full-credit description.” \\
    Medium & “Shows partial understanding; correctly defines terms but lacks complete comparison.” \\
    Low & “Mentions topic keywords without explanation; reasoning unclear and not linked to rubric.” \\

\end{longtable}
}

Overall, 87\% of explanations were rated high quality, 
9\% medium, and 4\% low.

%%%%%%%%%%%%%%%%%%%%%%%%%%%%%%%%%%%%%%%%%%%%%%%%%%%%%%%%%%%%

\section{LLM Scoring Error Categorization}

LLM scoring discrepancies were categorized into four main 
% error types. \rtab{tab:error_examples} summarizes their 
frequencies and provides representative cases.

{
\setlength{\LTleft}{-2.2cm}
\setlength{\LTright}{0pt}
\centering
\begin{longtable}{l c p{0.5\textwidth} l}
    \caption{Error Type Distribution and Illustrative Examples}
    \label{tab:error_examples}
    \\
    
    % For Header
    \hline
    \textbf{Error Type} & \textbf{Frequency (\%)} & \textbf{Description / Example} & \textbf{Bloom’s Level} \\
    \hline
    \endfirsthead

    % For Overflow Header
    \hline
    \textbf{Error Type} & \textbf{Frequency (\%)} & \textbf{Description / Example} & \textbf{Bloom’s Level} \\
    \hline
    \endhead

    % For Footer
    \hline
    \endfoot

    Under-scoring & 5.8 & AI rated 2 instead of 3 due to missing explicit keyword, despite correct reasoning. & Creating \\
    Over-scoring & 6.3 & AI gave 3 for an answer with factual inaccuracy but confident phrasing. & Analyzing \\
    Partial Credit Inconsistency & 3.4 & Similar responses received different scores due to varied sentence length. & Understanding \\
    Rubric Misinterpretation & 2.1 & Model misread rubric cue; rewarded example listing instead of explanation. & Applying \\
 
\end{longtable}
}


